\documentclass{article}
\usepackage[UTF8]{ctex}
\usepackage{amsmath,amssymb,amsfonts}
\usepackage{enumitem}
\usepackage{geometry}
\geometry{a4paper, margin=1in}
% === 错题本排版优化（不参与题目标识） ===
\setlength{\parindent}{0pt}
\setlength{\parskip}{0.35em plus 0.10em minus 0.08em}
\setlength{\emergencystretch}{3em}
\setlist{itemsep=0.25em, topsep=0.40em, parsep=0pt}
\allowdisplaybreaks[2]
% === 错题本排版优化结束 ===
\begin{document}
\section*{武汉大学2024年高等代数真题}
\section*{121．武汉大学2024年高等代数试卷}

\subsection{题目1}
一．（20 分）设数域 \(K\) 上的 \(n\) 阶矩阵 \(A\) 的第 \(i\) 行第 \(j\) 列的元素是 \(a_{i}-b_{j}\) 。

\begin{itemize}
  \item[（1）] 求行列式 \(|A|\) 。
  \item[（2）] \(n \geq 2, a_{1} \neq a_{2}, b_{1} \neq b_{2}\) 时，求方程组 \(A X=0\) 的解空间的维数和一个基．
\end{itemize}

二（20分）设 \(n\) 阶实方阵 \(A\) 的特征值全是实数，且 \(A\) 的所有一阶主子式之和与二阶主子式之和都是零。证明：\(A^{n}=O\) 。

\subsection{题目2}
三。（15 分）设 \(n\) 维欧氏空间 \(V\) 上的线性变换 \(\mathcal{T}\) 满足 \(\mathcal{T}^{3}+\mathcal{T}=\mathcal{O}\) ，证明： \(\mathcal{T}\) 的迹（即 \(\mathcal{T}\) 在 \(V\) 的某一基下对应矩阵的迹）等于零。

\subsection{题目3}
四．（20 分）设 \(A, B\) 是 \(n\) 阶实矩阵，且 \(A^{2}+B^{2}=A B\) ．证明：若 \(A B-B A\) 是可逆矩阵，则 \(n\) 是 3 的倍数。

\subsection{题目4}
五．设 \(A, B\) 是 \(n\) 阶复矩阵．证明：若 \(A B=B A\) ，那么存在 \(n\) 阶可逆复矩阵 \(P\) ，使得 \(P^{-1} A P, P^{-1} B P\)都是上三角矩阵。

\subsection{题目5}
六．（20分）设方阵 \(A\) 的特征值全是1．证明：对每个自然数 \(s, A^{s}\) 与 \(A\) 相似．

\subsection{题目6}
七．（20 分）设 \(f_{1}(x), f_{2}(x), \cdots, f_{n}(x)\) 为多项式 \((n \geq 2)\) ，且满足

\[
x^{n-1}+x^{n-2}+\cdots+x+1 \mid f_{1}\left(x^{n}\right)+x f_{2}\left(x^{n}\right)+\cdots+x^{n-1} f_{n}\left(x^{n}\right) .
\]

证明：存在常数 \(c\) ，使得 \((x-1)^{n} \mid \prod_{i=1}^{n}\left(f_{i}(x)-c\right)\) ．

\subsection{题目7}
八．（15 分）设 \(f(x), g(x)\) 是数域 \(K\) 上互素的多项式，\(A\) 是数域 \(K\) 上的 \(n\) 阶矩阵．证明：\(f(A) g(A)=\) \(O\) 的充分必要条件是 \(\operatorname{rank}(f(A))+\operatorname{rank}(g(A))=n\) ，其中 \(\operatorname{rank}\) 表示矩阵的秩。
\end{document}
